
\documentclass{article}
\usepackage{colortbl}
\usepackage{makecell}
\usepackage{multirow}
\usepackage{supertabular}

\begin{document}

\newcounter{utterance}

\centering \large Interaction Transcript for game `codenames', experiment `risk\_low', episode 2 with Apertus{-}v1.5{-}8B.
\vspace{24pt}

{ \footnotesize  \setcounter{utterance}{1}
\setlength{\tabcolsep}{0pt}
\begin{supertabular}{c@{$\;$}|p{.15\linewidth}@{}p{.15\linewidth}p{.15\linewidth}p{.15\linewidth}p{.15\linewidth}p{.15\linewidth}}
    \# & $\;$A & \multicolumn{4}{c}{Game Master} & $\;\:$B\\
    \hline

    \theutterance \stepcounter{utterance}  
    & & & \multicolumn{2}{p{0.3\linewidth}}{
        \cellcolor[rgb]{0.9,0.9,0.9}{
            \makecell[{{p{\linewidth}}}]{
                \texttt{\tiny{[GM$|$GM]}}
                {'hidden': {'team': ['aluniţă', 'lămâie', 'mașină', 'furculiţă', 'roman', 'tokyo', 'ketchup', 'perie', 'cangur'], 'innocent': ['baterie', 'africa', 'fructe de pădure', 'amazon', 'noapte', 'pirat', 'piept', 'sunet'], 'opponent': ['etichetă', 'film', 'val', 'alunecare', 'arc', 'zi', 'măr', 'mexic'], 'assassin': []}, 'revealed': {'team': {'team': [], 'innocent': [], 'opponent': [], 'assassin': []}, 'opponent': {'team': [], 'innocent': [], 'opponent': [], 'assassin': []}}}
            }
        }
    }
    & & \\ \\

    \theutterance \stepcounter{utterance}  
    & & \multicolumn{4}{p{0.6\linewidth}}{
        \cellcolor[rgb]{0.9,0.9,0.9}{
            \makecell[{{p{\linewidth}}}]{
                \texttt{\tiny{[P1$\langle$GM]}}
                \texttt{Furnizați un indiciu dintr{-}un singur cuvânt care se referă la unele dintre cuvintele echipei dvs. din lista de mai jos. Puteți alege la câte și la ce cuvinte se referă indiciul dumneavoastră.} \\
\texttt{Scopul este ca partenerul tău să ghicească toate cuvintele echipei tale pe baza indiciilor pe care i le dai, înainte ca cealaltă echipă să le descopere pe toate. Dacă partenerul dvs. ghicește un cuvânt "adversar" sau un cuvânt "distractor", rândul său de ghicit se încheie, iar echipa adversă are ocazia să facă un tur. Dacă partenerul tău ghicește un cuvânt "ucigaș", jocul se încheie și echipa ta pierde imediat.} \\
\texttt{Cu cât indiciul tău se referă la mai multe cuvinte, cu atât câștigi mai repede jocul, dar cu atât asocierile pot fi mai greu de ghicit pentru partenerul tău, așa că alege{-}ți cu înțelepciune cuvântul indiciu. Cuvântul indiciu trebuie să fie legat semantic de cuvintele țintă, nu poate fi unul dintre cuvintele din liste sau să conțină părți ale acestora.} \\
\\ 
\texttt{Furnizați întotdeauna indiciul cu un singur cuvânt și lista de cuvinte țintă, separate prin virgulă, în următorul format și faceți răspunsurile cât mai scurte posibil; nu includeți niciodată alt text decât este necesar în acest formular:} \\
\\ 
\texttt{INDICIU: <CUVÂNT>} \\
\texttt{ȚINTE: <ȚINTE>} \\
\\ 
\texttt{Cuvintele echipei tale sunt: aluniţă, lămâie, mașină, furculiţă, roman, tokyo, ketchup, perie, cangur.} \\
\texttt{Cuvintele echipei adversarului dvs. sunt: etichetă, film, val, alunecare, arc, zi, măr, mexic.} \\
\texttt{Cuvintele de distragere a atenției sunt: baterie, africa, fructe de pădure, amazon, noapte, pirat, piept, sunet.} \\
\texttt{Cuvintele de ucidere sunt: .} \\
            }
        }
    }
    & \\ \\

    \theutterance \stepcounter{utterance}  
    & \multicolumn{4}{p{0.6\linewidth}}{
        \cellcolor[rgb]{0.9,0.9,0.9}{
            \makecell[{{p{\linewidth}}}]{
                \texttt{\tiny{[P1$\rangle$GM]}}
                \texttt{INDICIU: formă} \\
\texttt{ȚINTE: aluniță, lămâie, măr} \\
            }
        }
    }
    & & \\ \\

    \theutterance \stepcounter{utterance}  
    & & & \multicolumn{2}{p{0.3\linewidth}}{
        \cellcolor[rgb]{0.9,0.9,0.9}{
            \makecell[{{p{\linewidth}}}]{
                \texttt{\tiny{[GM$|$GM]}}
                {'player': 'Cluegiver', 'type': 'target is hallucination', 'utterance': 'INDICIU: formă\nȚINTE: aluniță, lămâie, măr', 'message': "Targeted word 'aluniță' was not listed, you can only target words provided in the lists.", 'target': 'aluniță', 'previous_targets': [], 'remaining_words': ['lămâie', 'val', 'roman', 'baterie', 'amazon', 'piept', 'noapte', 'zi', 'măr', 'film', 'cangur', 'etichetă', 'furculiţă', 'aluniţă', 'sunet', 'mexic', 'africa', 'mașină', 'alunecare', 'fructe de pădure', 'pirat', 'ketchup', 'tokyo', 'perie', 'arc']}
            }
        }
    }
    & & \\ \\

\end{supertabular}
}

\end{document}
